\section{Threat Model and Security Statement}
\label{sec:threat-model}

We consider a prover that trains or serves a GBDT model for a verifier. The
prover may be malicious: it may commit to an arbitrary
(ill-formed) forest or inference output, and it may try to convince the
verifier that these objects arise from the public quantized training and
inference algorithms. The verifier is probabilistic polynomial time and checks
only the commitments, public parameters, public inputs, and proof transcript.

The protocol is organized into three phases.
\begin{itemize}
    \item \textbf{Committing phase.} The prover commits to the training dataset
    and the trained forest.

    \item \textbf{Forest phase.} The prover proves that the committed forest is
    well formed and, when training certification is requested, that it is
    consistent with a valid execution of the public GBDT trainer on the
    committed dataset.

    \item \textbf{Inference phase.} Given a batch of public inputs, the prover
    proves that the claimed inference output is consistent with the committed
    forest.
\end{itemize}

In the proof of training, the public statement consists of the parameters that
fix the training procedure, including the number of rounds $T$, the number of
classes $K$, the dataset size $M$, the feature dimension $d$, and the number of bins
$B$. The (optionally private) witness consists of the committed
dataset, labels, forest, and auxiliary training trace. The forest includes the
split metadata, tree topology, leaf boxes, and leaf values of all trees. Based
on the application, the dataset and the forest may be private or public.

\begin{theorem}[Informal security statement]
\label{thm:main-informal}
Assume the KZG commitment scheme is binding,
our proof system satisfies:
\begin{enumerate}
    \item \textbf{Completeness.} If the committed forest is obtained by honestly
    running the public quantized trainer on the committed dataset, then the
    prover can produce an accepting proof of training.
    If the committed forest is well-formed,
    the prover can produce an accepting proof of forest well-formedness and proof of inference for any public input batch.

    \item \textbf{Soundness.} If the verifier accepts, then with overwhelming
    probability, the committed forest is consistent with a valid execution of
    that same trainer on the committed dataset. Similarly, an accepting
    well-formedness proof certifies that the committed forest is well-formed, and
    an accepting
    inference proof certifies that the claimed output is obtained by evaluating
    the committed forest on the public input batch.

    \item \textbf{(Optional) Witness privacy.} 
    The proof system reveals nothing beyond the public statement. In particular, the proof of training hides 
    the training dataset, labels, forest, and auxiliary training trace, and the proof of well-formedness and inference hide the committed forest.
\end{enumerate}
\end{theorem}
